\exercisegroup

\begin{enumerate}
\def\labelenumi{\arabic{enumi})}
\tightlist
\item
  Let A be a closed and discrete subset of \(\mathbf{R}^2\) .
\end{enumerate}

\begin{enumerate}
\def\labelenumi{\alph{enumi})}
\item
  For R \textgreater{} 0, let \(D_{R}\) be the set of \(z \in C\) such that \(|z| < R\) ; show that the fundamental group of \(D_{R} - (A \cap D_{R})\) is isomorphic to the free group on \(A \cap D_{R}\) .
\item
  Prove that there exists a homeomorphism \(f\) of \(\mathbf{R}^{2}\) into itself such that \(f(\mathrm{A})\subset\mathbf{N}\times\{0\}\) .
\item
  Show that the fundamental group of \(\mathbf{R}^2 - \mathbf{A}\) is isomorphic to \(\mathrm{F(A)}\) .
\end{enumerate}

\begin{enumerate}
\def\labelenumi{\arabic{enumi})}
\setcounter{enumi}{1}
\item
  Let A be the set of points of \(R^{2}\) with coordinates \((1/n,0)\) , for \(n\in N^{*}\) . Show that the fundamental group of \(R^{2}-A\) has the cardinality of the continuum. In particular, it is not isomorphic to the group F(A).
\item
  Let \(((\mathrm{X}_{i}, x_{i}))_{i \in \mathrm{I}}\) be a family of pointed topological spaces; denote by \((\mathrm{X}, x)\) the wedge of this family (IV, p. 421, example 4). Suppose that, for every \(i \in I\) , the point \(x_{i}\) is closed in \(X_{i}\) and possesses a fundamental system of neighborhoods V such that the group \(\pi_{1}(\mathrm{V}, x_{i})\) is reduced to the identity element. Then, the canonical homomorphism \(*_{i \in \mathrm{I}} \pi_{1}(\mathrm{X}_{i}, x_{i}) \to \pi_{1}(\mathrm{X}, x)\) is an isomorphism.
\item
  Let P be the subspace of \(R^{2}\) defined as in exercise 5, I, p. 146, b). Let a be the point \((0,0,1)\) of \(R^{3}\) and let C be the union of the segments \([a,(x,0)]\) , for \(x \in P\) .
\end{enumerate}

\begin{enumerate}
\def\labelenumi{\alph{enumi})}
\item
  Show that the space C is locally arcwise connected and simply arcwise connected.
\item
  Let D be the union of the space C and its symmetric image -C. Show that the space D is simply connected, but that it is not simply arcwise connected.
\item
  Show that the Poincaré group \(\pi_{1}(D,0)\) has the cardinality of the continuum.
\end{enumerate}

\begin{enumerate}
\def\labelenumi{\arabic{enumi})}
\setcounter{enumi}{4}
\tightlist
\item
  Let P be the subspace of \(R^{2}\) defined as in exercise 5, I, p. 146, b); also denote by \(P^{+}\) (resp. \(P^{-}\) ) the set of points \((x,y)\in\mathrm{P}\) such that \(y\geqslant0\) (resp. \(y\leqslant0\) ). Let a be the point \((0,0,1)\) of \(R^{3}\) and let A be the union of the segment [a,0] and the segments \([a,(2r_{n},0)]\) , for \(n\in N\) . We set \(B^{+}=(P^{+}\times\{0\})\cup A\) , \(B^{-}=(P^{-}\times\{0\})\cup A\) and \(B=(P\times\{0\})\cup A\) .
\end{enumerate}

\begin{enumerate}
\def\labelenumi{\alph{enumi})}
\item
  Prove that the space \(B^{+}\) is loopable and that the Poincaré group \(\pi_{1}(B^{+}, a)\) is a free group on N.
\item
  Show that A is contractible to a.
\item
  Show that the canonical homomorphism from \(\pi_{1}(B^{+},a)*\pi_{1}(B^{-},a)\) into \(\pi_{1}(B,a)\) , deduced from the canonical inclusions of \(B^{+}\) and \(B^{-}\) into B, is not surjective.
\end{enumerate}

\begin{enumerate}
\def\labelenumi{\arabic{enumi})}
\setcounter{enumi}{5}
\tightlist
\item
  Let X be the real line R and let A be a subset of X. Let Y be the space deduced from X by contracting A to a point.
\end{enumerate}

\begin{enumerate}
\def\labelenumi{\alph{enumi})}
\item
  If A is discrete, the Poincaré group of Y is a free group.
\item
  If A has an accumulation point, the Poincaré group of Y has the cardinality of the continuum and is not free.
\item
  If A is the closure of the set of points \(1/n\) , for \(n \in \mathbf{N}^*\) , the space X/A is homeomorphic to the Hawaiian earring.
\end{enumerate}

\begin{enumerate}
\def\labelenumi{\arabic{enumi})}
\setcounter{enumi}{6}
\tightlist
\item
  Let G be a graph, S the set of its vertices and A the set of its oriented edges. We denote by \(|G|\) its geometric realization; we identify S with a subset of \(|G|\) and denote by \(p: I \times A \rightarrow |G|\) the canonical map.
\end{enumerate}

\begin{enumerate}
\def\labelenumi{\alph{enumi})}
\item
  Prove that the space \(|G|\) is compact (resp. locally compact) if and only if the graph G is finite (resp. locally finite).
\item
  Show that the canonical map from S to \(|\mathbf{G}|\) induces a bijection from \(\pi_0(\mathbf{G})\) onto \(\pi_0(|\mathbf{G}|)\) .
\item
  Denote by \(\varpi(|\mathrm{G}|)_\mathrm{S}\) the full subgroupoid of \(\varpi(|\mathrm{G}|)\) with vertex set S. Show that there exists a unique isomorphism of groupoids \(\varpi (\mathrm{G})\to \varpi (|\mathrm{G}|)_{\mathrm{S}}\) which is the identity on the set of vertices and which maps the class of an oriented edge \(a\in \mathrm{A}\) to the class of the path \(t\mapsto p(t,a)\)
\item
  Suppose that the graph G is connected and finite; let \(m = 1 + \text{Card}(S) - \text{Card}(A)\) . Show that the Poincaré group of \(|G|\) is a free group with m generators.
\end{enumerate}

\begin{enumerate}
\def\labelenumi{\arabic{enumi})}
\setcounter{enumi}{7}
\tightlist
\item
  Let X be the quotient space of the space \(R \times \{0,1\}\) by the coarsest equivalence relation identifying \((x,0)\) and \((x,1)\) for all \(x \in R^{*}\) (“real line with doubled origin”).
\end{enumerate}

\begin{enumerate}
\def\labelenumi{\alph{enumi})}
\item
  Prove that \(\pi_1(\mathrm{X},x)\) is isomorphic to \(\mathbf{Z}\) , for every point \(x\in \mathrm{X}\) .
\item
  Let \(\mathrm{X}^{(2)}\) be the quotient space of the space \(X^{2}\) by the action of the symmetric group \(S_{2}\) on the space \(X^{2}\) acting by permutation of coordinates. Show that \(\pi_{1}(\mathrm{X}^{(2)},y)=0\) for all \(y\in\mathrm{X}^{(2)}\) .
\end{enumerate}

\begin{enumerate}
\def\labelenumi{\arabic{enumi})}
\setcounter{enumi}{8}
\tightlist
\item
  Let X be a topological space, let R be an equivalence relation on X and let \(p: X \to X/R\) be the canonical surjection. Let \(x \in X\) .
\end{enumerate}

\begin{enumerate}
\def\labelenumi{\alph{enumi})}
\item
  Suppose that X is connected and locally arcwise connected, that the equivalence classes are connected sets, and that the space X/R is loopable. Show that the homomorphism \(p_{*}:\pi_{1}(X,x)\to\pi_{1}(X/R,p(x))\) is surjective.
\item
  Suppose that X = I and that \(\{0,1\}\) is the only equivalence class which is not reduced to a single element. Show that the homomorphism \(p_{*}\) is not surjective.
\end{enumerate}

\begin{enumerate}
\def\labelenumi{\arabic{enumi})}
\setcounter{enumi}{9}
\tightlist
\item
  Let W be the Warsaw circle, defined in exercise 1, IV, p. 455, whose notations we take up again.
\end{enumerate}

\begin{enumerate}
\def\labelenumi{\alph{enumi})}
\item
  Let R be the equivalence relation on W whose only non-singleton classes are B and S; let \(p: W \rightarrow W/R\) be the canonical surjection. Show that the quotient space W/R is homotopy equivalent, but not homeomorphic, to the circle \(S_{1}\) . Show that the homomorphism \(p_{*}\) is not surjective.
\item
  Let R' be the equivalence relation on W for which B ∪ S is the only equivalence class not reduced to a single element; let \(p'\) : W → W/R' be the canonical surjection. Show that the space W/R' is homeomorphic to \(S_{1}\) , that the homomorphism \(p_{*}\) is not surjective, but that the equivalence classes of R' are connected.
\end{enumerate}

\begin{enumerate}
\def\labelenumi{\arabic{enumi})}
\setcounter{enumi}{10}
\tightlist
\item
  Let X be a path-connected topological space and let \(f: X \to X\) be a homeomorphism. Let T be the quotient space of \(X \times I\) by the finest equivalence relation for which the points \((x,0)\) and \((f(x),1)\) are equivalent, for every \(x \in X\) . We denote \(p: X \times I \to T\) the canonical map; we denote by \(j\colon X\to T\) the map given by \(x\mapsto p(x,1/2)\) . Let a be a point of X, let \(c\in\Lambda_{f(a),a}(X)\) be a path of origin \(f(a)\) and of term \(a\) , and let \(\gamma=[c]\) .
\end{enumerate}

\begin{enumerate}
\def\labelenumi{\alph{enumi})}
\item
  Show that the map from I to T given by \(t \mapsto p(a, \frac{1}{2} - t)\) for \(t \in [0, \frac{1}{2}]\) and \(t \mapsto p(c(2t - 1), \frac{3}{2} - t)\) for \(t \in ]\frac{1}{2}, 1]\) is a loop at \(j(a)\) in T. We denote \(\delta\) its class.
\item
  Let S be the unique group homomorphism from \(\pi_{1}(\mathrm{X},a)*\mathbf{Z}\) into \(\pi_{1}(\mathrm{T},j(a))\) which maps a class \(v\in\pi_{1}(\mathrm{X},a)\) to \(j_{*}(v)\) and the element t=1 of Z to the class \(\delta\). Show that the homomorphism S is surjective and that its kernel is the smallest normal subgroup containing the elements \(vt(\gamma^{-1}f_{*}(v)\gamma)^{-1}t^{-1}\) , for \(v\in\pi_{1}(\mathrm{X},a)\) .
\item
  For \(v \in \pi_{1}(X, a)\) we set \(\varphi(v) = \gamma^{-1} f_{*}(v) \gamma\) ; prove that \(\varphi\) is an automorphism of the group \(\pi_{1}(X, a)\) . Let \(\varphi: \mathbf{Z} \to \operatorname{Aut}(\pi_{1}(X, a))\) be the unique group homomorphism that maps 1 to \(\varphi\) . Construct an isomorphism from the group \(\pi_{1}(T, j(a))\) onto the external semidirect product of Z by \(\pi_{1}(X, a)\) with respect to \(\varphi\) .
\end{enumerate}

\begin{enumerate}
\def\labelenumi{\arabic{enumi})}
\setcounter{enumi}{11}
\tightlist
\item
  \begin{enumerate}
  \def\labelenumii{\alph{enumii})}
  \tightlist
  \item
    Let D be a line in \(R^{3}\) and let \(U = R^{3} - D\) . Prove that U is path-connected and that for every point \(a \in U\) , the group \(\pi_{1}(U, a)\) is isomorphic to Z.
  \end{enumerate}
\end{enumerate}

\begin{enumerate}
\def\labelenumi{\alph{enumi})}
\setcounter{enumi}{1}
\tightlist
\item
  Let P be a plane in \(R^{3}\), let C be a circle contained in P and let \(V = R^{3} - C\) . Prove that V is path-connected and that, for every point \(a \in V\) , the group \(\pi_{1}(V, a)\) is isomorphic to Z. (Apply the van Kampen theorem to the closed covering of \(R^{3}\) formed by the two half-spaces bounded by P; one may also note that the spaces U and V are homeomorphic.)
\end{enumerate}

\begin{enumerate}
\def\labelenumi{\arabic{enumi})}
\setcounter{enumi}{12}
\tightlist
\item
   Let D and D' be two distinct lines of \(R^{3}\), and let \(U = R^{3} - (D \cup D')\) .
\end{enumerate}

\begin{enumerate}
\def\labelenumi{\alph{enumi})}
\item
  Prove that U is path-connected.
\item
  Assume that D and D' are disjoint; prove that, for every point \(a \in U\) , the group \(\pi_{1}(U, a)\) is a free group on two generators.
\item
  Assume that D and D' have a common point; prove that, for every point \(a \in U\) , the group \(\pi_{1}(U, a)\) is a free group on three generators.
\end{enumerate}

\begin{enumerate}
\def\labelenumi{\arabic{enumi})}
\setcounter{enumi}{13}
\tightlist
\item
   Let P be a plane in \(R^{3}\), let C be a circle contained in P and let D be a line in \(R^{3}\) . Let \(U = R^{3} - (C \cup D)\) and let a be a point of U.
\end{enumerate}

\begin{enumerate}
\def\labelenumi{\alph{enumi})}
\item
  Prove that U is path-connected.
\item
  Assume that D does not intersect the convex hull of C. Prove that \(\pi_{1}(U,a)\) is a free group on two generators.
\item
  Assume that D does not meet C but does meet the convex hull of C. Prove that \(\pi_{1}(U,a)\) is isomorphic to \(Z^{2}\) .
\item
  Assume that D intersects C at exactly one point. Prove that \(\pi_{1}(U,a)\) is a free group on two generators.
\item
  Assume that D intersects C at two points. Prove that \(\pi_{1}(U,a)\) is a free group on three generators.
\end{enumerate}

\begin{enumerate}
\def\labelenumi{\arabic{enumi})}
\setcounter{enumi}{14}
\tightlist
\item
  A topological space X is said to have the Phragmén–Brouwer property if it is connected and if, for every pair (A, B) of disjoint closed subsets of X such that \(X-A\) and \(X-B\) are connected, then \(X-(A \cup B)\) is connected.
\end{enumerate}

\begin{enumerate}
\def\labelenumi{\alph{enumi})}
\item
  Let X be a locally connected topological space which has the Phragmén–Brouwer property. Let (A, B) be a pair of disjoint closed subsets of X, and let \((a, b)\) be a pair of points of \(X-(A \cup B)\) which belong to the same connected component of \(X-A\) (resp. of \(X-B\)). Prove that a and b belong to the same connected component of \(X-(A \cup B)\) .
\item
  Let X be a connected and locally connected by arcs topological space that does not have the Phragmén–Brouwer property. Prove that, for every point \(x \in X\), there exists a surjective homomorphism from \(\pi_{1}(X, x)\) into Z.
\item
  Let X be a connected and locally connected by arcs topological space that is simply connected by arcs. Prove that X has the Phragmén-Brouwer property.
\end{enumerate}

d) Let X be a separated and locally connected topological space such that, for every point \(x \in \mathrm{X}\) , the space \(\mathrm{X} - \{x\}\) has the Phragmén-Brouwer property. Let A be a subset of X that is homeomorphic to \([0,1]\) ; prove that \(\mathrm{X} - \mathrm{A}\) is connected. (Let \(c: \mathbf{I} \to \mathrm{A}\) be a homeomorphism. Argue by contradiction and construct a sequence \((\mathrm{J}_n)\) where, for every \(n\) , \(\mathrm{J}_n\) is an interval of \(\mathbf{I}\) of length \(2^{-n}\) , contained in \(\mathrm{J}_{n-1}\) if \(n \geqslant 1\) , such that \(\mathrm{X} - c(\mathrm{J}_n)\) is not connected; let \(x\) be the unique point of X common to all the sets \(c(\mathrm{J}_n)\) Then show that \(\mathrm{X} - \{x\}\) is not connected.)

\begin{enumerate}
\def\labelenumi{\alph{enumi})}
\setcounter{enumi}{4}
\tightlist
\item
  Let \(n\) be an integer \(\geqslant 1\) and let A be a subset of \(\mathbf{S}_n\) which is homeomorphic to \([0,1]\) . Prove that \(\mathbf{S}_n - \mathbf{A}\) is connected.
\end{enumerate}

\begin{enumerate}
\def\labelenumi{\arabic{enumi})}
\setcounter{enumi}{15}
\tightlist
\item
  Let A be a subset of \(S_{2}\) which is homeomorphic to \(S_{1}\) . Let a and b be distinct points of A; we denote by \(A'\) and \(A''\) the connected components of \(A-\{a,b\}\) , and one sets \(X=S_{2}-\{a,b\}\) , \(U=X-A'\) and \(V=X-A''\) .
\end{enumerate}

\begin{enumerate}
\def\labelenumi{\alph{enumi})}
\item
  Show that U and V are connected. (Use Exercise 15.)
\item
  Show that for every point \(x \in U\) the homomorphism from \(\pi_{1}(U, x)\) into \(\pi_{1}(X, x)\) deduced from the inclusion of U in X is trivial.
\item
  Prove that \(S_{2}-A\) has exactly two connected components and that their boundary is equal to A.
\item
  Let S be a subset of \(R^{2}\) which is homeomorphic to \(S_{1}\) ; prove that \(R^{2}-S\) has exactly two connected components and that their boundary is equal to S (Jordan's theorem).
\end{enumerate}

\begin{enumerate}
\def\labelenumi{\arabic{enumi})}
\setcounter{enumi}{16}
\tightlist
\item
  We identify \(S_{1}\) with the set of complex numbers of modulus 1, and the sphere \(S_{3}\) with the subspace of \(C^{2}\) consisting of the points \((z,w)\) such that \(|z|^{2} + |w|^{2} = 1\) . Let B and C be the subspaces of \(S_{3}\) defined by \(|z| \leqslant |w|\) and \(|z| \geqslant |w|\) respectively.
\end{enumerate}

\begin{enumerate}
\def\labelenumi{\alph{enumi})}
\item
  Show that B and C are homeomorphic to \(B_{2} \times S_{1}\) .
\item
  For every pair \((m,n)\) of relatively prime integers, we denote by \(K_{m,n}\) the image of the map from \(S_{1}\) into \(S_{3}\) given by \(u\mapsto(u^{m},u^{n})\) ("torus knot").
\item
  Prove that \(S_{3}-K_{m,n}\) is connected and its fundamental group admits a presentation \(\langle x,y;x^{m}=y^{n}\rangle\) .
\end{enumerate}

d) Prove that there exists a homeomorphism \(\varphi\) of \(\mathbf{S}_3\) into itself such that \(\varphi(\mathrm{K}_{m,n}) = \mathrm{K}_{1,0}\) if and only if \(|m| \leqslant 1\) or \(|n| \leqslant 1\) .

\begin{enumerate}
\def\labelenumi{\arabic{enumi})}
\setcounter{enumi}{17}
\tightlist
\item
  Let \((m,n)\) be a pair of natural integers such that 1 \textless{} m \textless{} n. Let G be the group defined by the presentation \(\langle x,y;x^{m}=y^{n}\rangle\) . Let \(z=x^{m}\) and let C be the subgroup of G generated by z.
\end{enumerate}

\begin{enumerate}
\def\labelenumi{\alph{enumi})}
\item
  Prove that C is contained in the center of G.
\item
  Prove that G/C is isomorphic to the free product \((\mathbf{Z}/m\mathbf{Z})*(\mathbf{Z}/n\mathbf{Z})\) .
\item
  Deduce that C is the center of G.
\item
  Let \((p,q)\) be a pair of natural integers such that 1 \textless{} p \textless{} q. Assume that the group \(G'\) defined by the presentation \(\langle x,y; x^{p} = y^{q}\rangle\) is isomorphic to G. Prove that \((p,q) = (m,n)\) .
\item
  Show that the quotient of G by its derived subgroup is isomorphic to Z.
\end{enumerate}

\begin{enumerate}
\def\labelenumi{\arabic{enumi})}
\setcounter{enumi}{18}
\tightlist
\item
  Let G be a graph; denote by S the set of its vertices and by A the set of its directed edges, and \(|\mathbf{G}|\) its geometric realization. Denote \(j\colon \mathrm{S}\to |\mathrm{G}|\) and \(p\colon \mathbf{I}\times \mathbf{A}\to |\mathbf{G}|\) the canonical maps. We assume that the sets S and A are finite.
\end{enumerate}

 A map \(f\colon|G|\to\mathbf{R}^{n}\) is said to be affine (resp. piecewise affine) if for every edge \(a\in A\) , the map \(t\mapsto f\circ p(a,t)\) is affine (resp. piecewise affine). (Cf. III, p. 331, exerc. 6.)

\begin{enumerate}
\def\labelenumi{\alph{enumi})}
\item
   Prove that there exists an injective and piecewise affine map from \(|G|\) into \(R^{3}\) (Begin by showing that there exists an integer \(n \geqslant 1\) and an injective, piecewise affine map from \(|G|\) into \(R^{n}\) .)
\item
   For every continuous injective map \(f\colon|G|\to R^{n}\), prove that there exists a continuous injective map \(g\colon|G|\to R^{n}\) which is piecewise affine and homotopic to f.
\item
   We assume that the map \((o,t)\colon\mathrm{A}\to\mathrm{S}\times\mathrm{S}\) is injective. Let f be an injective map from S to R. Prove that there exists a unique affine map F from \(|G|\) into \(R^{3}\) which maps every vertex \(s\in S\) onto the point \((f(s),f(s)^{2},f(s)^{3})\) of \(R^{3}\) . Prove that F is injective.
\end{enumerate}

We say that the graph G is planar if there exists a continuous injective map from \(|G|\) into \(R^{2}\) .

d) We assume that the graph \(\mathbf{G}\) is planar and that the map \((o,t)\colon\mathbf{A}\to\mathbf{S}\times\mathbf{S}\) is injective. Let \(f\colon|\mathbf{G}|\to\mathbf{R}^{2}\) be a continuous injective map. Prove that there exists an injective affine map \(g\colon|\mathbf{G}|\to\mathbf{R}^{2}\) which is homotopic to \(f\) .

\begin{enumerate}
\def\labelenumi{\arabic{enumi})}
\setcounter{enumi}{19}
\tightlist
\item
   Let \(f \colon \mathbf{I} \to \mathbf{R}^2\) be a piecewise affine loop whose restriction to [0,1[ is injective; denote \(C = f(\mathbf{I})\) .
\end{enumerate}

\begin{enumerate}
\def\labelenumi{\alph{enumi})}
\item
   Prove that \(\mathbf{R}^2 -\mathbf{C}\) has at most two connected components.
\item
   For \(x \in R^{2} - C\) and \(v \in S_{1}\) , let \(n_{v}(x)\) be the number of connected components of \((\mathbf{R}^{2} - \mathbf{C}) \cap (x + \mathbf{R}_{+}v)\) . Prove that there exists a unique locally constant map \(n: R^{2} - C \to \{0,1\}\) such that \(n(x) \equiv n_{v}(x) \pmod{2}\) for all \(v \in S_{1}\) .
\item
  Deduce that \(R^{2}-C\) has exactly two connected components and that their boundaries are equal to C.
\item
   Let \(g: I \to R^{2}\) be a piecewise affine injective map such that \(\overline{g}^{-1}(C) = \{0,1\}\) ; we set \(P = g(I)\) and \(D = C \cup P\) . Prove that \(C - (C \cap P)\) has two connected components; denote them by \(P_{1}\) and \(P_{2}\) . Prove that \(R^{2} - D\) has three connected components and that their boundaries are equal to C, \(P_{1} \cup P\) and \(P_{2} \cup P\) .
\item
   Let G be a finite planar graph; let S be the set of its vertices and A the set of its oriented edges. Let \(f\colon|G|\to\mathbf{R}^{2}\) be a continuous injective piecewise affine map and let P be its image. Prove that the number of connected components of \(\mathbf{R}^{2}-P\) is equal to \(1+\mathrm{Card(S)}-\mathrm{Card(A)}+\mathrm{Card}(\pi_{0}(\mathrm{G}))\) .
\end{enumerate}

\begin{enumerate}
\def\labelenumi{\arabic{enumi})}
\setcounter{enumi}{20}
\item
  Let K be the graph associated with the quiver whose set of vertices is \(\{1,2,3\} \times \{0,1\}\) , the set of arrows is the set of pairs of the form \(((a,0),(b,1))\) , for \(a,b\in \{1,2,3\}\) and whose source and target maps are respectively derived from the first and second projection (“Kuratowski graph”). Prove that the graph K is not planar.
\item
  Let \(n\) be a natural integer and let \(\mathrm{K}_n\) be a complete graph whose vertex set has cardinality \(n\) (II, p. 220, exerc. 5).
\end{enumerate}

\begin{enumerate}
\def\labelenumi{\alph{enumi})}
\item
  Prove that the graphs \(K_{n}\) , for \(1 \leqslant n \leqslant 4\) , are planar.
\item
  Prove that the graph \(K_{5}\) is not planar.
\end{enumerate}

\begin{enumerate}
\def\labelenumi{\arabic{enumi})}
\setcounter{enumi}{22}
\tightlist
\item
  Let \(u\) and \(v\) be two points of \(\mathbf{S}_1\) such that \(\mathcal{I}(v / u) > 0\) ; one defines a family \((a_j)_{j \in \mathbf{Z} / 4\mathbf{Z}}\) by \(a_1 = u\) , \(a_2 = v\) , \(a_3 = -u\) and \(a_4 = -v\) . Denote by \(\mathrm{Q}_1\) and \(\mathrm{Q}_2\) the two segments \([a_1, a_3]\) and \([a_2, a_4]\) in the unit disk \(\mathbf{B}_2\) and let us set \(\mathrm{Q} = \mathrm{Q}_1 \cup \mathrm{Q}_2\) .
\end{enumerate}

\begin{enumerate}
\def\labelenumi{\alph{enumi})}
\tightlist
\item
  Prove that, for every \(j \in Z/4Z\) , there exists a unique arcwise connected component \(A_{j} \in \pi_{0}(B_{2}-Q)\) whose closure contains \(a_{j}\) and \(a_{j-1}\) . Prove that the union of the family \((A_{j})\) is equal to \(B_{2}-Q\) .
\end{enumerate}

 Let us set \(C = B_{2} \times [-1, 1]\) , \(N = (0, 0, 1)\) and \(S = (0, 0, -1)\) . Let \(f: [-1, 1] \to R_{+}\) be a continuous map such that \(f(-1) = f(1) = 0\) and \(0 < f(t) < 1\) for all \(t \in ]-1; 1[\) . Let us then denote by \(L_{1}\) and \(L_{2}\) the images of the maps from \([-1, 1]\) into C given by \(t \mapsto (ta_{1}, 0)\) and \(t \mapsto (ta_{2}, f(t))\) respectively; finally set \(L = L_{1} \cup L_{2}\) .

\begin{enumerate}
\def\labelenumi{\alph{enumi})}
\setcounter{enumi}{1}
\item
  Prove that C-L is arcwise connected.
\item
  We say that a path \(c\) in \(\mathbf{C} - \mathbf{L}\) with origin \(\mathbf{N}\) and endpoint \(\mathbf{S}\) is straight if there exist juxtaposed paths \(c_1\) and \(c_2\) in \(\mathbf{B}_2\) such that \(c\) is the juxtaposition of the paths given by \(t \mapsto (c_1(t), 1)\) , \(t \mapsto (c_1(1), 1 - 2t)\) and \(t \mapsto (c_2(t), -1)\) . Prove that then \(c_1(1) \not\in \mathbb{Q}\) ; we denote by \(\mathrm{A}_c\) the path-connected component of \(\mathbf{B}_2 - \mathbf{Q}\) which contains this point.
\end{enumerate}

d) Let \(c\) and \(c'\) be paths in \(\mathrm{C} - \mathrm{L}\), with origin \(\mathrm{N}\) and endpoint \(\mathrm{S}\) ; suppose they are straight. Prove that they are strictly homotopic if and only if \(\mathrm{A}_c = \mathrm{A}_{c'}\) .

\begin{enumerate}
\def\labelenumi{\alph{enumi})}
\setcounter{enumi}{4}
\tightlist
\item
  For every \(j \in \mathbf{Z}/4\mathbf{Z}\) , let \(c_j\) be a path in \(C-L\) with origin \(N\) and endpoint \(S\) which is straight and such that \(A_{c_j} = A_j\) . Prove the relation
\end{enumerate}

\[
[ c _ {1} ] [ c _ {4} ] ^ {- 1} = [ c _ {2} ] [ c _ {3} ] ^ {- 1}
\]

in the group \(\pi_1(\mathrm{C} - \mathrm{L},\mathrm{N})\)

\begin{enumerate}
\def\labelenumi{\alph{enumi})}
\setcounter{enumi}{5}
\tightlist
\item
   Prove that the canonical homomorphism from the free group on two generators \(x, y\) into \(\pi_1(\mathrm{C} - \mathrm{L}, \mathrm{N})\) which maps \(x\) and \(y\) to \([c_1][c_2]^{-1}\) and \([c_2][c_3]^{-1}\) respectively is an isomorphism of groups. (Apply the van Kampen theorem to the covering of \(\mathrm{C} - \mathrm{L}\) formed by the eight closed sets delimited by the three vector planes of \(\mathbf{R}^3\) containing two of the points of the set \(\{a_1, a_2, \mathrm{N}\}\) .)
\end{enumerate}

\begin{enumerate}
\def\labelenumi{\arabic{enumi})}
\setcounter{enumi}{23}
\tightlist
\item
   Let m be an integer \(\geqslant 1\) and let \((\theta_{j})_{1\leqslant j\leqslant m}\) be a sequence of real numbers such that \(0\leqslant\theta_{1}<\theta_{2}<\cdots<\theta_{n}<2\pi\) . For \(j\in Z/mZ\) we set \(a_{j}=\exp(i\theta_{k})\) , where k is the unique element of \(\{1,\ldots,n\}\) belonging to class j; we also denote by \(Q_{j}\) the segment \([0,a_{j}]\) of \(B_{2}\) . Finally we set \(Q=Q_{1}\cup\cdots\cup Q_{m}\) .
\end{enumerate}

\begin{enumerate}
\def\labelenumi{\alph{enumi})}
\item
  For \(j \in Z/mZ\) , prove that there exists a unique path-connected component \(A_{j}\) of \(B_{2}-Q\) which contains every point \(b \in S_{1}\) such that \(\mathcal{I}(b/a_{j}) > 0\) and \(\mathcal{I}(b/a_{j+1}) < 0\) . Show that the family \((A_{j})\) is formed by pairwise disjoint sets and that its union is equal to \(B_{2}-Q\) .
\item
   Set \(C = B_{2} \times [-1, 1]\) , \(N = (0, 0, 1)\) , \(S = (0, 0, -1)\) , and \(L = Q \times \{0\}\) . We say that a path in \(C-L\), with origin N and endpoint S, is straight if there exist paths \(c_{1}\) and \(c_{2}\) in \(B_{2}\) , juxtaposed, such that c is the juxtaposition of the paths given by \(t \mapsto (c_{1}(t), 1)\) , \(t \mapsto (c_{1}(1), 1 - 2t)\) and \(t \mapsto (c_{2}(t), -1)\) . Prove that then \(c_{1}(1) \notin Q\) ; we denote by \(A_{c}\) the path-connected component of \(B_{2}-Q\) which contains this point.
\item
   Let c and \(c'\) be paths in \(C-L\), with origin N and endpoint S; suppose they are straight. Prove that they are strictly homotopic if and only if \(A_{c} = A_{c'}\) .
\end{enumerate}

d) For every \(j \in \mathbf{Z}/m\mathbf{Z}\) , let \(c_j\) be a loop in \(C-L\), with origin N and endpoint S, which is straight and such that \(\mathrm{A}_{c_j} = \mathrm{A}_j\) . Prove that the homomorphism from the free group \(\mathrm{F}(\mathbf{Z}/m\mathbf{Z})\) into \(\pi_1(\mathrm{C}-\mathrm{L},\mathrm{N})\) which maps \(x_j\) to \([c_j][c_{j+1}]^{-1}\) , for \(j \in \mathbf{Z}/m\mathbf{Z}\) , is surjective, and that its kernel is the smallest normal subgroup which contains the product \(x_1x_2 \ldots x_m\) .

\begin{enumerate}
\def\labelenumi{\arabic{enumi})}
\setcounter{enumi}{24}
\tightlist
\item
   Set \(C = B_{2} \times [-1, 1]\) . Let K be a closed subspace of \(R^{3}\) and let P be its image under the projection given by \((x, y, z) \mapsto (x, y)\) . We suppose that there exist a finite set I, a family \((p_{i})_{i \in \mathrm{I}}\) of points of \(R^{2}\) , a family \((r_{i})_{i \in \mathrm{I}}\) of strictly positive real numbers and a family \((m_{i})_{i \in \mathrm{I}}\) of natural numbers satisfying the following properties, where \(\varphi_{i}\) denotes the map from \(R^{3}\) into \(R^{3}\) given by \(x \mapsto (p_{i}, 0) + r_{i}x\) and \(C_{i} = \varphi_{i}(C) \) :
\end{enumerate}

- The sets \(C_{i}\) are pairwise disjoint;

 - For every \(i \in \mathrm{I}\) the set \(\mathrm{L}_i = \overline{\varphi}_i^{-1}(\mathrm{C}_i \cap \mathrm{K})\) is the subspace \(\mathrm{L}\) defined in exercise 23 if \(m_i = 0\) and in exercise 24 (where we set \(m = m_i\) ) if \(m_i \geqslant 1\) .

- The set \(\mathrm{K}' = \mathrm{K} - \bigcup_{i} (\mathring{\mathrm{C}}_i \cap \mathrm{K})\) is the union of a finite family of closed segments pairwise disjoint in the plane \(\mathbf{R}^2 \times \{0\}\) .

Let \(\mathrm{N} = (x_{\mathrm{N}}, y_{\mathrm{N}}, z_{\mathrm{N}})\) and \(\mathrm{S} = (x_{\mathrm{S}}, y_{\mathrm{S}}, z_{\mathrm{S}})\) be points of \(R^{3}\) such that \(z_{N} > \sup(r_{i})_{i \in I}\) and \(z_{\mathrm{S}} < -\sup(r_{i})_{i \in I}\) .

\begin{enumerate}
\def\labelenumi{\alph{enumi})}
\tightlist
\item
   We say that a path \(c\) in \(\mathbf{R}^3-K\) with origin N and endpoint S is straight if there exist paths \(c_1\) and \(c_2\) in \(\mathbf{R}^2\) , juxtaposed, such that \(c\) is the juxtaposition of the paths given by \(t \mapsto (c_1(t), z_{\mathrm{N}})\) , \(t \mapsto (c_1(1), (1 - t)z_{\mathrm{N}} + tz_{\mathrm{S}})\) , and \(t \mapsto (c_2(t), z_{\mathrm{S}})\) .
\end{enumerate}

Prove that then the point \(c_1(1)\) does not belong to P; we denote by \(A_c\) the path-connected component of \(\mathbf{R}^2 - \mathbf{P}\) which contains this point. Prove that two straight paths \(c\) and \(c'\) are strictly homotopic if and only if \(A_c = A_{c'}\) .

   For every path-connected component \(\alpha\) of \(\mathbf{R}^2-\mathbf{P}\) , we fix a path \(c(\alpha)\) in \(\mathbf{R}^2-\mathbf{K}\) , with origin N and endpoint S which is straight and such that \(\mathrm{A}_{c(\alpha)} = \alpha\) .

\begin{enumerate}
\def\labelenumi{\alph{enumi})}
\setcounter{enumi}{1}
\item
   Let \(i \in I\) be such that \(m_{i} = 0\) and such that \(\overline{\varphi}_{i}^{-1}(K \cap C_{i})\) is the set L defined in exercise 23, whose notations we take up again; we have \(\overline{\varphi}_{i}^{-1}(P \times \{0\} \cap C_{i}) = Q \times \{0\}\) . For all \(j \in Z/4Z\) , denote by \(\alpha_{i,j}\) the path-connected component of \(R^{2} - P\) which contains the image under \(\varphi_{i}\) of the connected component \(A_{j}\) of \(B_{2} - P\) . Show that we have \([c(\alpha_{i,1})][c(\alpha_{i,4})]^{-1} = [c(\alpha_{i,2})][c(\alpha_{i,3})]^{-1}\) in \(\pi_{1}(R^{3} - K, N)\) .
\item
   Let \(\omega\) be an element of \(\pi_0(\mathbf{R}^2-\mathbf{P})\) . Let \(\lambda_\omega\colon \mathrm{F}(\pi_0(\mathbf{R}^2-\mathbf{P}))\to \pi_1(\mathbf{R}^3-\mathbf{K},\mathbf{N})\) be the unique group homomorphism which, for \(\alpha \in \pi_0(\mathbf{R}^2-\mathbf{P})\) , maps the element \(x_\alpha\) onto \([c(\alpha)][c(\omega)]^{-1}\) .
\end{enumerate}

Prove that the homomorphism \(\lambda_{\omega}\) is surjective and its kernel is the smallest normal subgroup of \(\mathrm{F}(\pi_0(\mathbf{R}^2 - \mathbf{P}))\) which contains the element \(x_{\omega}\) and the elements \(x_{\alpha_{i,1}}^{-1}x_{\alpha_{i,2}}x_{\alpha_{i,3}}^{-1}x_{\alpha_{i,4}}\) , for \(i \in \mathrm{I}\) such that \(m_i = 0\) .

\begin{enumerate}
\def\labelenumi{\alph{enumi})}
\setcounter{enumi}{3}
\item
   For every segment \(s\) belonging to \(\pi_{0}(\mathrm{K}^{\prime})\) , we choose elements \(\alpha_{s}\) and \(\alpha_{s}^{\prime}\) of \(\pi_{0}(\mathbf{R}^{2}-\mathbf{P})\) such that the set of connected components of \(R^{2}-P\) whose closure contains \(s\) is equal to \(\{\alpha_{s},\alpha_{s}^{\prime}\}\) . Let \(\mu\colon\mathrm{F}(\pi_{0}(\mathrm{K}^{\prime}))\to\pi_{1}(\mathbf{R}^{3}-\mathbf{K},\mathbf{N})\) be the unique group homomorphism which, for \(s\in\pi_{0}(\mathrm{K}^{\prime})\) , maps the element \(x_{s}\) onto \([c(\alpha_{s})][c(\alpha_{s}^{\prime})]^{-1}\) . Prove that \(\mu\) is surjective. For \(i\in I\) such that \(m_{i}=0\) and \(j\in\{1,2,3,4\}\) , let \(s_{i,j}\) be the unique segment of \(\pi_{0}(\mathrm{K}^{\prime})\) such that \(\varphi_{i}(a_{j})\) meets \(s_{i,j}\) . Prove that there exist elements \(u_{i},v_{i},w_{i}\) of \(\{-1,1\}\) such that \(\mu(x_{s_{i,4}})=\mu(x_{s_{i,2}}^{w_{i}})\) and \(\mu(x_{s_{i,3}})=\mu(x_{s_{i,2}}^{-u_{i}}x_{s_{i,1}}^{v_{i}}x_{s_{i,2}}^{u_{i}})\) . Show that the kernel of \(\mu\) is the smallest normal subgroup of \(\mathrm{F}(\pi_{0}(\mathrm{K}^{\prime}))\) which contains the elements \(x_{s_{i,3}}x_{s_{i,2}}^{-u_{i}}x_{s_{i,1}}^{-v_{i}}x_{s_{i,2}}^{u_{i}}\) and \(x_{s_{i,4}}x_{s_{i,2}}^{-w_{i}}\) for \(i\in I\) .
\item
  Set
\end{enumerate}

\[
k = \operatorname{Card} (\pi_ {0} (\mathrm{K})) - \operatorname{Card} (\{i \in \mathrm{I} \mid m _ {i} > 0 \}) + \frac {1}{2} \sum_ {i \in \mathrm{I}} m _ {i}.
\]

   Prove that \(k \in N\) . Prove that the quotient of \(\pi_{1}(\mathbf{R}^{3} - \mathbf{K})\) by its derived subgroup is isomorphic to \(Z^{k}\) .

\begin{enumerate}
\def\labelenumi{\arabic{enumi})}
\setcounter{enumi}{25}
\tightlist
\item
  In the numerical plane \(\mathbf{R}^2\) , let \(\mathrm{D}_0\) be the disk with center \((0,0)\) and radius 1 and \(\mathrm{C}_0\) its boundary; we set \(a_0 = (-1,0)\) . Let \(g\) be a natural integer \(\geqslant 1\) , let \(r\) be a strictly positive real number and let \((v_1,\ldots,v_g)\) be a sequence of real numbers such that \(r < \inf (v_{1} + 1,(v_{2} - v_{1}) / 2,\dots,(v_{g} - v_{g - 1}) / 2,1 - v_{g})\) . For all \(j\in \{1,\ldots,g\}\) , we denote by \(\mathrm{D}_j\) the closed disk with center \((0,v_j)\)
\end{enumerate}

and of radius \(r\) and by \(\mathrm{C}_j\) the circle with center \((0,v_j)\) and of radius \(r\) ; we also set \(a_{j}=(-r,v_{j})\) . Let X be the space \(\mathrm{D}=\bigcup_{j=1}^{g}\mathring{\mathrm{D}}_{j}\) .

  For every \(j \in \{1, \ldots, j\}\) , let \(u_j\) be the class in \(\varpi(\mathrm{X})\) of a path with origin \(a_0\) and of term \(a_j\) whose image is the segment \([a_0, a_j]\) ; we also denote by \(c_j \in \pi_1(\mathrm{X}, a_j)\) the class of the loop given by \(t \mapsto (-r \cos(2\pi t), -r \sin(2\pi t) + v_j)\) and \(\gamma_j = u_j c_j u_j^{-1}\) .

\begin{enumerate}
\def\labelenumi{\alph{enumi})}
\item
  Let \(\gamma_{0}\) be the class of the loop in X given by \(t \mapsto (-\cos(2\pi t), -\sin(2\pi t))\) . Show that we have \(\gamma_{0} = \gamma_{1} \ldots \gamma_{g}\) .
\item
  Prove that the unique homomorphism from the free group on g generators \(x_{1},\ldots,x_{j}\) into \(\pi_{1}(X,a_{0})\) which maps \(x_{j}\) to \(\gamma_{j}\) is a group isomorphism. (Apply the van Kampen theorem to the covering of X defined by its intersections with the half-planes \(R_{+}\times R\) and \(R_{-}\times R\) .)
\item
  Let \(f \colon \mathrm{X} \to \mathbf{R}_+\) be a continuous function such that \(f^{-1}(0) = \bigcup_{j=0}^{g} \mathrm{C}_j\) and let Y be the subspace of \(\mathbf{R}^3\) consisting of the points \((x, y, z)\) such that \(z^2 = f(x, y)\) . We denote \(i_+\) (resp. \(i_-\) ) the continuous map from X to Y given by \((x, y) \mapsto (x, y, \sqrt{f(x, y)})\) (resp. \((x, y) \mapsto (x, y, -\sqrt{f(x, y)})\) ). For every path class \(\gamma\) in X, one defines path classes in \(\varpi(\mathrm{Y})\) by \(\gamma^+ = (i_+)_{*}(\gamma)\) and \(\gamma^- = (i_-)_*(\gamma)\) .
\end{enumerate}

For every \(j \in \{1, \ldots, g\}\) , one finally defines an element of \(\pi_1(Y, a_0)\) by \(\delta_j = u_j^+ u_j^-\) . Prove that one has the relations \(\gamma_j^+ \delta_j = \delta_j \gamma_j^-\) for \(1 \leqslant j \leqslant g\) , and \(\gamma_1^+ \ldots \gamma_g^+ = \gamma_1^- \ldots \gamma_g^-\) .

\begin{enumerate}
\def\labelenumi{\alph{enumi})}
\setcounter{enumi}{3}
\tightlist
\item
  Let \(\varphi\) be the unique homomorphism from a free group F with 2g generators \(x_{1},\ldots ,x_{g},y_{1},\ldots ,y_{g}\) into \(\pi_1(\mathrm{Y},a_0)\) which maps \(x_{j}\) to \(\gamma_j^+\) and \(y_{j}\) to \(\delta_{j}\) . Show that the homomorphism \(\varphi\) is surjective and its kernel is the smallest normal subgroup of F that contains the element
\end{enumerate}

\[
(x _ {1} ^ {- 1} y _ {1} x _ {1}) (x _ {2} ^ {- 1} y _ {2} x _ {2}) \dots (x _ {g} ^ {- 1} y _ {g} x _ {g}) (y _ {1} \dots y _ {g}) ^ {- 1}.
\]

(Apply the van Kampen theorem to the covering of Y defined by its intersections with the half-spaces \(R^{2} \times R_{+}\) and \(R^{2} \times R_{-}\) .)

\begin{enumerate}
\def\labelenumi{\alph{enumi})}
\setcounter{enumi}{4}
\tightlist
\item
  Let \(\varphi'\) be the unique homomorphism from a free group F with 2g generators \(x_1, \ldots, x_g, y_1, \ldots, y_g\) into \(\pi_1(Y, a_0)\) which, for every \(j\) , maps \(x_j\) to
\end{enumerate}

\[
(\gamma_ {1} ^ {+} \dots \gamma_ {j - 1} ^ {+}) ^ {- 1} \delta_ {j} (\gamma_ {1} ^ {+} \dots \gamma_ {j - 1} ^ {+})
\]

 and \(y_{j}\) to \(\gamma_{j}^{+}\) . Show that the homomorphism \(\varphi'\) is surjective and its kernel is the smallest normal subgroup of F that contains the element

\[
\left(x _ {1} ^ {- 1} y _ {1} ^ {- 1} x _ {1} y _ {1}\right) \dots \left(x _ {g} ^ {- 1} y _ {g} ^ {- 1} x _ {g} y _ {g}\right).
\]

\begin{enumerate}
\def\labelenumi{\arabic{enumi})}
\setcounter{enumi}{26}
\tightlist
\item
  Let D be the unit disk in C and let X be a subspace of D which is a neighborhood of \(S_{1}\) . Let R be the coarsest equivalence relation in X which, for \(u, v \in S_{1}\) , identifies the points \(f(u)\) and \(f(v)\) if \(f(u) = f(v)\) . Let \(p\colon X \to X / R\) be the canonical surjection. Let \(T\) be a compact topological space and let \(f\colon S_1 \to T\) be a continuous surjective map. We set \(a = f(1)\) and denote by \(\alpha \in \pi_1(T, a)\) the class of the loop \(t \mapsto f(e^{2\pi it})\) in \(T\) .
\end{enumerate}

\begin{enumerate}
\def\labelenumi{\alph{enumi})}
\item
  Prove that there exists a unique map \(\sigma\) from T to X/R such that \(\sigma \circ f = p|_{\mathbf{S}_1}\) . Prove that \(\sigma\) defines a homeomorphism from T onto its image \(p(\mathbf{S}_1)\) .
\item
  Let r be a real number such that 0 \textless{} r \textless{} 1; we suppose that X is the set of \(z \in D\) such that \(r \leqslant |z| \leqslant 1\) . Show that X/R is homeomorphic to the mapping cylinder of f. Deduce that the homomorphism \(\sigma_{*}\) is a group isomorphism from \(\pi_{1}(T, a)\) onto \(\pi_{1}(X/R, p(1))\) .
\item
  Suppose that X = D. Show that X/R is homeomorphic to the mapping cone of f. Deduce that the group homomorphism \(\sigma_{*}\colon\pi_{1}(\mathrm{T},a)\to\pi_{1}(\mathrm{X}/\mathrm{R},p(1))\) is surjective and its kernel is the smallest normal subgroup of \(\pi_{1}(\mathrm{T},a)\) which contains \(\alpha\) .
\item
  Let r be a real number such that 0 \textless{} r \textless{} 1; let \(X_{1}\) be the set of \(z \in X\) such that \(|z| \leqslant r\) and \(X_{2}\) the set of \(z \in D\) such that \(r \leqslant |z| \leqslant 1\) ; we suppose that \(X_{2} \subset X\) . Let \(\beta \in \pi_{1}(X_{2}, r)\) be the class of the loop given by \(t \mapsto re^{2\pi it}\) ; let \(\delta \in \varpi_{p(r), p(1)}\) be the class of the path given by \(t \mapsto p((1 - t)r + t)\) .
\end{enumerate}

  Prove that there exists a unique group homomorphism \(\mu\) from \(\pi_1(\mathrm{T},a)*\pi_1(\mathrm{X}_2,r)\) into \(\pi_1(\mathrm{X} / \mathrm{R},p(1))\) which coincides with the homomorphism \(\sigma_*\) on \(\pi_1(\mathrm{T},a)\) and with the homomorphism \(v\mapsto \delta^{-1}p_{*}(v)\delta\) on \(\pi_1(\mathrm{X}_2,p(r))\) . Prove that \(\mu\) is surjective and its kernel is the smallest normal subgroup containing \(\alpha \beta^{-1}\) .

\begin{enumerate}
\def\labelenumi{\arabic{enumi})}
\setcounter{enumi}{27}
\tightlist
\item
  Let X be a topological space. A homotopy \(\sigma: \mathrm{X} \times \mathbf{I} \to \mathrm{X}\) is said to be an isotopy if, for every \(t \in \mathbf{I}\) , the map \(x \mapsto \sigma(x, t)\) deduced from \(\sigma\) by passage to the subspace is a homeomorphism of X onto itself.
\end{enumerate}

\begin{enumerate}
\def\labelenumi{\alph{enumi})}
\item
  Let m be a natural integer and let \(f: X \to R^{n}\) be a continuous map. Prove that there exists an isotopy whose origin is the identity map of \(X \times R^{n}\) and whose term is the map from \(X \times R^{n}\) into \(X \times R^{n}\) given by \((x, y) \mapsto (x, y - f(x))\) .
\item
  Let A and B be subspaces of X; one says that A is isotopic to B if there exists an isotopy \(\sigma\colon\mathbf{X}\times\mathbf{I}\to\mathbf{X}\) whose origin is the identity map on X and such that \(f(\mathbf{A}\times\{1\})=\mathbf{B}\) .
\end{enumerate}

Prove that the relation “A is isotopic to B” is an equivalence relation on the set of subspaces of X (resp. on the set of closed subspaces of X).

\begin{enumerate}
\def\labelenumi{\alph{enumi})}
\setcounter{enumi}{2}
\tightlist
\item
  Let m and n be natural numbers, and let A be a closed subspace of \(R^{m}\) and let B be a closed subspace of \(R^{n}\) . Assume that A and B are homeomorphic. Prove that the subspaces \(A \times \{0\}\) and \(\{0\} \times B\) of \(R^{m+n}\) are isotopic.
\end{enumerate}

\begin{enumerate}
\def\labelenumi{\arabic{enumi})}
\setcounter{enumi}{28}
\tightlist
\item
  Let A be a closed subset of \(R^{2}\) homeomorphic to a closed subset B of R.
\end{enumerate}

\begin{enumerate}
\def\labelenumi{\alph{enumi})}
\item
  Prove that the subspaces \(A \times \{0\}\) and \(\{(0,0)\} \times B\) of \(R^{3}\) are isotopic.
\item
  We assume that \(B \neq R\) ; prove that \(R^{2}-A\) is path-connected.
\item
  Assume that B = R; prove that \(R^{2}-A\) has exactly two path-connected components and that A is the boundary of each of them.
\item
  Let A be a subset of \(S_{2}\) homeomorphic to \(S_{1}\) ; prove that \(S_{2} - A\) has exactly two path-connected components and that A is the boundary of each of them.
\item
  Let A be a subset of \(R^{2}\) homeomorphic to \(S_{1}\) ; prove that \(R^{2}-A\) has exactly two path-connected components and that A is the boundary of each of them.
\end{enumerate}

\begin{enumerate}
\def\labelenumi{\arabic{enumi})}
\setcounter{enumi}{29}
\tightlist
\item
  Let \(v = (v_{1},\ldots ,v_{n})\colon \mathbf{R}^{n}\to \mathbf{R}^{n}\) be a map of class \(\mathrm{C}^1\) and let \(\delta\) be a real number \(>0\) ; we make the following hypotheses:
\end{enumerate}

- The support of \(v\) is contained in \([- \delta, 1 + \delta] \times \mathring{\mathbf{B}}_{n-1}\) ;

- For every \(x \in [0, 1 - \delta[\) , we have \(v_1(x, 0, \ldots, 0) > 0\) , and \(v_1(1, 0, \ldots, 0) = 0\) ;

- For every \(x \in [0,1]\) , we have \(v_{2}(x,0,\ldots,0) = \cdots = v_{n}(x,0,\ldots,0) = 0\) .

Let \(\Phi\colon\mathbf{R}^{n}\times\mathbf{R}\to\mathbf{R}^{n}\) be the integral flow of the map \((x,t)\mapsto v(x)\) (VAR, 9.1.3); for \(t\in\mathbf{R}\) , let \(\Phi_{t}\) be the map \(x\mapsto\Phi(x,t)\) .

\begin{enumerate}
\def\labelenumi{\alph{enumi})}
\item
  Prove that there exists a real number \(\tau\) such that \(\Phi_{\tau}([0,1] \times \{0\}) \subset [1 - \delta, 1] \times \{0\}\) .
\item
  Let \(j\colon[-\delta,1+\delta]\times\mathbf{B}_{n-1}\to\mathbf{R}^{n}\) be a continuous injective map. For every \(t\in[0,1]\) we set \(\mathrm{A}_{t}=j([t,1]\times\{0\})\) . Prove that, for every \(t\in]0,1[\) , the subspaces \(A_{0}\) and \(A_{t}\) are isotopic.
\item
  Let P be a subset of \(\mathbf{R}^{n}\) homeomorphic to \([0,1]\) which is the union of a finite family of segments. Prove that P is isotopic to a segment.
\end{enumerate}

\begin{enumerate}
\def\labelenumi{\arabic{enumi})}
\setcounter{enumi}{30}
\tightlist
\item
  Let \(f = (f_1, \ldots, f_n) \colon \mathbf{I} \to \mathbf{R}^n\) be a map of class \(\mathbf{C}^1\) such that \(f'(t) \neq 0\) for all \(t \in \mathbf{I}\) .
\end{enumerate}

\begin{enumerate}
\def\labelenumi{\alph{enumi})}
\item
  We assume that \(f_{1}^{\prime}(t) > 0\) for all \(t \in I\) . For all \(j \in \{2, \ldots, n\}\) , prove that there exists a continuous map \(g_{j} \colon R \to R\) such that \(f_{j}(x) = g_{j}(f_{1}(x))\) for all \(x \in I\) . Deduce that \(f(\mathbf{I})\) is isotopic to a segment.
\item
   Prove that there exists a real number \(\delta > 0\) and a continuous injective map \(j\colon [-\delta, 1 + \delta] \times \mathbf{B}_{n-1} \to \mathbf{R}^n\) such that \(j(t, 0) = f(t)\) for all \(t \in \mathbf{I} .\) Deduce from Exercise 30 that, for every \(\tau \in [0, 1[\) the set \(f(\mathbf{I})\) is isotopic to \(f([\tau, 1])\) .
\item
  Prove that \(f(\mathbf{I})\) is isotopic to a segment.
\end{enumerate}

\begin{enumerate}
\def\labelenumi{\arabic{enumi})}
\setcounter{enumi}{31}
\tightlist
\item
  Let A be a closed subset of \(\mathbf{R}^3\) homeomorphic to \([0,1]\) . The endpoints of A are the two points \(p\) of A such that \(\mathrm{A} - \{p\}\) is connected (cf. III, p. 280, th. 2); the other points of A will be called interior.
\end{enumerate}

Let \(p\) be a point of A; one says that A is tame at \(p\) if there exists a neighborhood V of \(p\) in \(\mathbf{R}^3\) and a homeomorphism from V onto the unit ball of \(\mathbf{R}^3\) which maps \(\mathrm{A} \cap \mathrm{V}\) onto a segment.

\begin{enumerate}
\def\labelenumi{\alph{enumi})}
\item
  Let p be a point of A such that A is tame at p. Let \((\mathrm{V}_{n})\) be a decreasing sequence of neighborhoods of p in \(R^{3}\) such that \(\bigcap_{n}\mathrm{V}_{n}=\{p\}\) for all n; for every n, let \(a_{n}\) be a point of \(V_{n}-A\) and denote by \(j_{n}\colon\pi_{1}(\mathrm{V}_{n}-\mathrm{A},a)\to\pi_{1}(\mathrm{V}_{1}-\mathrm{A},a)\) the group homomorphism induced by the inclusion of \(V_{n}\) into \(V_{1}\). If p is an endpoint of A (resp. an interior point), prove that there exists an integer m such that for every integer \(n\geqslant m\) the image of the homomorphism \(j_{n}\) is reduced to the identity element (resp. is abelian).
\item
  Assume that A is tame at every point. Prove that there exists a homeomorphism from \(R^{3}\) onto itself that maps A onto a segment. Deduce that \(R^{3}-A\) is connected and simply connected by arcs.
\end{enumerate}

\begin{enumerate}
\def\labelenumi{\arabic{enumi})}
\setcounter{enumi}{32}
\tightlist
\item
  \begin{enumerate}
  \def\labelenumii{(\arabic{enumii})}
  \setcounter{enumii}{5}
  \tightlist
  \item
    Let L be a subspace of \(\mathbf{R}^3\) as defined in Exercise 23, where one takes \(u = (1 - i)/\sqrt{2}\) and \(v = (1 + i)/\sqrt{2}\) . Let \(a\) and \(r\) be real numbers such that \(0 < 2r < a < 1 - 2r\) . We set \(p_1 = (0, -a)\) , \(p_2 = (0, 0)\) and \(p_3 = (0, a)\) ; \(q_1 = (-1, -a)\) , \(q_2 = (-1, 0)\) , \(q_3 = (-1, a)\) ; \(q_1' = (1, -a)\) , \(q_2' = (1, 0)\) , \(q_3' = (1, a)\) . Let K be the union of the sets \(p_1 + rL\) , \(p_2 + rL\) , \(p_3 + rL\) and the segments \([q_1, p_1 - rv]\) , \([q_2, p_2 - ru]\) , \([q_3, p_3 - ru]\) , \([q_1', p_1 + ru]\) , \([q_2', p_3 + ru]\) , \([q_3', p_3 + rv]\) .
  \end{enumerate}
\end{enumerate}

Let \((x_{n})_{n\in \mathbf{Z}}\) be a strictly increasing family of real numbers such that \(\lim_{n\to -\infty}x_n = -1\) and \(\lim_{n\to +\infty}x_n = 1\) ; for \(n\in \mathbf{Z}\) we set \(z_{n} = y_{n} = 1 - |x_{n}|\) . For \(n\in \mathbf{Z}\) , let \(f_{n}\) be the map from \([-1,1]^{3}\) into \(\mathbf{R}^3\) given by

\[
f _ {n} (x, y, z) = \frac {1 - x}{2} (x _ {n}, y y _ {n}, z z _ {n}) + \frac {1 + x}{2} (x _ {n + 1}, y y _ {n + 1}, z z _ {n + 1})
\]

, and we set \(A_{n}=f_{n}(K)\) and \(B_{n}=f_{n}([-1,1]^{3})\) . Let A be the union of the family \((\mathrm{A}_{n})_{n\in\mathbf{Z}}\) .

\begin{enumerate}
\def\labelenumi{\alph{enumi})}
\item
  Prove that A is homeomorphic to [0, 1].

(6) This example is due to E. ARTIN and R. FOX, “Some wild cells and spheres in three-dimensional space,” Annals of Math. 49 (1948), pp. 979–990.

\item
  For every natural integer \(m\) we set
\end{enumerate}

\[
\mathrm{A} _ {m} ^ {\prime} = \bigcup_ {n <   - m} \mathrm{B} _ {n} \cup \bigcup_ {- m \leqslant n \leqslant m} \mathrm{A} _ {m} \cup \bigcup_ {n > m} \mathrm{B} _ {n}.
\]

Prove that \(R^{3}-A_{m}^{\prime}\) is path-connected and that its fundamental group is generated by elements \(a_{n}, b_{n}, c_{n}\) (for \(n \in Z\) such that \(-m \leqslant n \leqslant m\) ) subject to the relations \(b_{-m} = c_{-m}a_{-m}\) , \(c_{n+1}a_{n+1} = c_{n}c_{n+1}\) , \(c_{n+1}b_{n} = a_{n}c_{n+1}\) , \(b_{n}c_{n+1} = b_{n+1}b_{n}\) (for \(n \in Z\) such that \(-m \leqslant n < m\) ) and \(c_{m}a_{m} = b_{m}\) .

\begin{enumerate}
\def\labelenumi{\alph{enumi})}
\setcounter{enumi}{2}
\tightlist
\item
   Prove that \(R^{3}-A\) is path-connected and that its fundamental group is generated by elements \(a_{n}, b_{n}, c_{n}\) (for \(n \in Z\) ) subject to the relations
\end{enumerate}

\[
b _ {n} = c _ {n} a _ {n}, \quad c _ {n + 1} a _ {n + 1} = c _ {n} c _ {n + 1}, \quad c _ {n + 1} b _ {n} = a _ {n} c _ {n + 1}, \quad b _ {n} c _ {n + 1} = b _ {n + 1} b _ {n}
\]

 for \(n \in Z\) .

d) Prove that there exists a homomorphism from \(\pi_1(\mathbf{R}^3 -\mathbf{A})\) into the group \(\mathfrak{S}_5\) which maps the class of \(c_{n}\) to the permutation \((1,2,3,4,5)\mapsto (2,3,4,5,1)\) if \(n\) is odd and to the permutation \((1,2,3,4,5)\mapsto (4,3,5,2,1)\) if \(n\) is even. Deduce that \(\mathbf{R}^3 -\mathbf{A}\) is not simply connected by arcs.

\begin{enumerate}
\def\labelenumi{\arabic{enumi})}
\setcounter{enumi}{33}
\tightlist
\item
  Let X be a path-connected topological space, let G be a discrete group acting continuously on the left on X. Let M be an open subset of X such that G·M = X. Define the sets S, T, the group F, and the group homomorphism \(\varphi\colon\mathrm{F}\to\mathrm{G}\) as in prop. 10 of I, p. 136. Let \(a\) be a point of M.
\end{enumerate}

\begin{enumerate}
\def\labelenumi{\alph{enumi})}
\item
  Let \(c\colon\mathbf{I}\to\mathbf{X}\) be a loop at \(a\) . Prove that there exists a unique element \(g(c)\) of F satisfying the following property: for every integer \(n\) and every sequence \((s_{1},\ldots,s_{n})\) of elements of S such that \(c([(k-1)/n,k/n])\subset s_{k}\mathrm{M}\) for every integer \(k\in\{1,\ldots,n\}\) , we have \(g(c)=x_{s_{n}^{-1}s_{1}}x_{s_{1}^{-1}s_{2}}\ldots x_{s_{n-1}^{-1}s_{n}}\) .
\item
  Prove that there exists a unique group homomorphism \(\gamma\) of \(\pi_1(\mathrm{X},a)\) into F such that \(\gamma ([c]) = g(c)\) for every loop \(c\) in X at \(a\) .
\item
  Prove that the kernel of \(\varphi\) is equal to the image of the homomorphism \(\gamma\) .
\item
  Prove that the kernel of \(\gamma\) is generated by the union of the images of the groups \(\pi_{1}(\mathrm{M}\cup s\cdot\mathrm{M},a)\) , for \(s\in S\) .
\end{enumerate}

